\chapter{Συμπεράσματα και Μελλοντικές Επεκτάσεις}
\label{ch:conclusions}

Στο \cref{ch:results} τα αποτελέσματα αξιολογήθηκαν τόσο σε εξομοίωση (Spike ISS)
όσο και σε RISC-V hardware, και τοποθετήθηκαν έναντι υλοποιήσεων
διαφορετικών αρχιτεκτονικών κατηγοριών\ifTODO, ISA ablation, και constant-time
ασφάλειας\fi. Το παρόν κεφάλαιο
συνοψίζει τη συνεισφορά της εργασίας, αναγνωρίζει τους περιορισμούς της, και
χαρτογραφεί τις κατευθύνσεις μελλοντικής έρευνας.

%% ==========================================================================
\section{Σύνοψη Συνεισφοράς}
\label{sec:summary}

Η παρούσα εργασία παρουσίασε την πρώτη πλήρως βελτιστοποιημένη υλοποίηση του \gls{hqc} \gls{kem}
σε RISC-V σε καθαρό λογισμικό (χωρίς custom hardware), αξιοποιώντας αποκλειστικά επικυρωμένες επεκτάσεις bit-manipulation και vector-crypto (\gls{zbc},
\gls{zbb}, \gls{zvbc}, \gls{zvbb}). Η υλοποίηση οργανώθηκε σε τρία tiers:

\paragraph{Tier~1: Scalar (\Zbc/\Zbb)}
Ο πολυωνυμικός πολλαπλασιασμός επιταχύνθηκε κατά \speedup{101.3} (HQC-128) μέσω
Karatsuba compile-time specialized tree generation και εγγενούς \clmul/\clmulh. Το πλήρες
\gls{kem} επιτυγχάνει \speedup{35.0} στο HQC-128 και
\speedup{60.1} στο HQC-256.

\paragraph{Tier~2: Base-RVV (\RVV + \Zbc/\Zbb)}
Η προσθήκη του \RVV 1.0 αυξάνει το polymul
speedup σε \speedup{124.6} (HQC-128) και το πλήρες \gls{kem} σε \speedup{55.9},
με υβριδική αρχιτεκτονική (RVV polymul + scalar Keccak).

\paragraph{Tier~3: Vector-Crypto (\Zvbc/\Zvbb)}
Η πλήρης αξιοποίηση \gls{zvbc} \insn{vclmul.vx} επιτυγχάνει \speedup{264.1} στο
polymul (HQC-128) και \speedup{436.5} στο HQC-256. Το πλήρες \gls{kem} φτάνει
\speedup{75.9} (HQC-128) και \speedup{163.8} (HQC-256). Σε επίπεδο δυναμικών
εντολών (\instret) αυτό είναι
της ίδιας τάξης μεγέθους με τον x86-64 FFT+CRT champion. Σε εκτιμώμενους κύκλους
ρολογιού το χάσμα είναι μεγαλύτερο, και
τόσο η άμεση σύγκριση \instret-έναντι-κύκλων όσο και η επικύρωση σε πραγματικό
silicon παραμένουν μελλοντική εργασία (\cref{sec:future-work}).

\paragraph{Keccak και αποκωδικοποίηση}
Πέραν του πολυωνυμικού πολλαπλασιασμού, βελτιστοποιήθηκαν ο Keccak
(\speedup{1.37} μέσω scalar \gls{zbb} idiom recognition) και η
αποκωδικοποίηση \gls{rm}/\gls{rs} (\speedup{14.27} μέσω
\RVV vectorization και \gls{zvbc} $\GF{2^8}$ acceleration).

\ifTODO
\paragraph{Constant-time ασφάλεια}
Η timing side-channel ασφάλεια επαληθεύθηκε μέσω διπλής στρατηγικής:
αρχιτεκτονική ανάλυση Zkt (τεκμηρίωση ότι κάθε εντολή ανήκει στο
data-independent timing subset) και εμπειρική στατιστική ανίχνευση dudect
σε hardware (BPI-F3 / SpacemiT K1 X60) για τα Tier~1 και Tier~2
variants (\cref{sec:ct-verification}). Η εμπειρική dudect κάλυψη του Tier~2
περιορίζεται σε επίπεδο polymul. Η αρχιτεκτονική ανάλυση Zkt τεκμηριώνει
θεωρητική CT συμμόρφωση για το Tier~3.
\fi

\paragraph{Hardware validation}
Τα Tier~1 και Tier~2 επικυρώθηκαν σε πραγματικό RISC-V hardware (BPI-F3,
SpacemiT K1 X60, VLEN $= \num{256}$), παρέχοντας μετρήσεις κύκλων με pipeline
effects (\cref{sec:hw-evaluation}). Το scalar path επιτυγχάνει IPC $\approx 1.49$. Συνεπώς ο \instret δεν αποτελεί
απόλυτο cycle predictor (υπερεκτιμά τους κύκλους κατά περίπου \pct{50}), αλλά παρέχει
συνεπή, αναλογικό εκτιμητή κύκλων για in-order scalar κώδικα, όπου ο
σταθερός IPC δίνει κύκλους $\approx \instret/\text{IPC}$.

%% ==========================================================================
\section{Περιορισμοί}
\label{sec:limitations}

Η παρούσα εργασία υπόκειται στους ακόλουθους περιορισμούς, οι οποίοι
οριοθετούν την εμβέλεια των συμπερασμάτων.

\paragraph{Απουσία Tier~3 hardware validation}
Το Tier~3 (\gls{zvbc}/\gls{zvbb}) αξιολογήθηκε αποκλειστικά σε εξομοίωση (Spike ISS)
λόγω μη διαθεσιμότητας, κατά την εκπόνηση, εμπορικού silicon με \gls{zvbc} υποστήριξη. Το πρώτο τέτοιο SoC (SpacemiT K3) κυκλοφόρησε μόλις τον Μάιο 2026 και επομένως δεν ήταν διαθέσιμο κατά τη διεξαγωγή της εργασίας. Τα
αποτελέσματα αντιπροσωπεύουν \instret (αριθμό δυναμικά αποσυρθεισών εντολών). Η πραγματική απόδοση σε silicon αναμένεται να
διαφοροποιηθεί ανάλογα με τη μικροαρχιτεκτονική υλοποίηση του \gls{zvbc} execution unit.

\paragraph{Μικροαρχιτεκτονική ετερογένεια στη cross-architecture σύγκριση}
Οι συγκρινόμενες υλοποιήσεις εκτελούνται σε δομικά διαφορετικούς επεξεργαστές
(in-order έναντι out-of-order, dual- έναντι multi-issue, soft-cores σε FPGA), 
οπότε η άμεση αριθμητική σύγκριση απαιτεί
επιφυλάξεις. Ο \instret αποτελεί κάτω φράγμα κύκλων σε in-order pipelines, αλλά σε
superscalar/dual-issue πυρήνες (όπως ο SpacemiT X60 με IPC $\approx 1.49$) γίνεται
άνω φράγμα, και σε κάθε περίπτωση δεν ενσωματώνει ILP, OoO execution, ή cache
effects. Επιπλέον, η hardware validation των Tier~1/2 στηρίζεται σε έναν μόνο
in-order πυρήνα (SpacemiT K1 X60), οπότε η μεταφορά των cycle-based συμπερασμάτων
σε out-of-order silicon ή διαφορετικό vector pipeline μπορεί να μεταβάλει τον IPC.

\ifTODO
\paragraph{Εμπειρική CT verification μόνο για Tier~1/2}
Η dudect εμπειρική επαλήθευση καλύπτει τα variants \code{scalar\_ref}, \code{scalar\_zb}
(Tier~1) και \code{vector\_rvv} (Tier~2). Το \code{vector\_zvk}
(Tier~3) δεν δοκιμάστηκε με dudect λόγω απουσίας \gls{zvbc}/\gls{zvbb} στον SpacemiT K1.
Η αρχιτεκτονική ανάλυση Zkt (\cref{ssec:ct-design-patterns}) τεκμηριώνει
θεωρητική CT συμμόρφωση, αλλά η εμπειρική επιβεβαίωση σε πραγματικό silicon
παραμένει εκκρεμής.

\fi

\paragraph{Απουσία ενεργειακού χαρακτηρισμού}
Όλες οι μετρικές αφορούν performance proxies (\instret/κύκλοι). Δεν πραγματοποιήθηκαν
μετρήσεις energy-per-operation ή power envelope, μετρικές που θα ολοκλήρωναν την
εικόνα απόδοσης σε deployment σενάρια πέραν του instruction/cycle proxy.

%% ==========================================================================
\section{Μελλοντική Έρευνα}
\label{sec:future-work}

Η εργασία ανοίγει πολλαπλές κατευθύνσεις μελλοντικής έρευνας, οι οποίες
κατηγοριοποιούνται σε βραχυπρόθεσμες (εφικτές με υπάρχοντα εργαλεία) και
μακροπρόθεσμες (απαιτούν νέο hardware ή ISA extensions).

%% --------------------------------------------------------------------------
\ifTODO
\subsection{Constant-Time Υλοποίηση}
\label{sec:future-ct}

\paragraph{Formal CT verification για RVV/Zvbc binaries}
Η τυπική επαλήθευση ιδιοτήτων σταθερού χρόνου σε RVV/\gls{zvbc} binaries αποτελεί
ανοιχτό ερευνητικό πρόβλημα. Τα υπάρχοντα εργαλεία (Binsec/Rel, ct-verif,
HolBA) δεν υποστηρίζουν τη RISC-V vector ISA
(\cref{ssec:ct-formal-gap}). Η ανάπτυξη formal verification
framework για RVV θα επέτρεπε μαθηματική απόδειξη, αντί εμπειρικής
τεκμηρίωσης, της CT ιδιότητας.

\paragraph{Εμπειρική dudect validation του Tier~3}
Με τη διαθεσιμότητα RVA23-compliant hardware (\eg ο SpacemiT K3 ή μεταγενέστερα SoCs με πλήρη
\gls{zvbc}/\gls{zvbb}), η εκτέλεση dudect σε πραγματικό silicon θα παρέχει εμπειρική
επιβεβαίωση ότι η \gls{zvbc} pipeline δεν εισάγει data-dependent timing variations.
Η υπάρχουσα dudect υποδομή (\code{tests/ct\_dudect/}) είναι έτοιμη για
deployment σε νέο hardware χωρίς τροποποιήσεις.

\paragraph{Publish-mode dudect (10M measurements)}
Η τρέχουσα dudect εκτέλεση χρησιμοποιεί $\sim$\num{100000} measurements ανά
trial. Η επέκταση σε \num{10000000} measurements ανά trial (publish-mode) θα
ενισχύσει τη στατιστική ισχύ, μειώνοντας τον κίνδυνο false negatives σε
ασθενείς διαρροές.
\fi

%% --------------------------------------------------------------------------

\paragraph{Tier~3 performance validation σε πραγματικό silicon}
Η κεντρική εκκρεμότητα είναι η μέτρηση πραγματικών κύκλων για το
\code{vector\_zvk} σε hardware με εγγενή \gls{zvbc} υποστήριξη. Η μέτρηση
αυτή θα αποκαλύψει τον πραγματικό CPI του \insn{vclmul.vx} και θα
επιβεβαιώσει (ή θα αναθεωρήσει) τα \instret-based speedup claims.

\paragraph{Cycle-accurate RTL simulation (Marian/Ara soft-core)}
Η χρήση του Marian ή Ara open-source RISC-V vector processor σε FPGA ή
Verilator θα παρέχει χαρακτηρισμό της pipeline με ακρίβεια κύκλου για \gls{zvbc}
εντολές, γεφυρώνοντας το χάσμα μεταξύ Spike \instret και πραγματικών κύκλων
πριν τη διαθεσιμότητα commercial silicon.

%% --------------------------------------------------------------------------

\paragraph{Vectorization εκκρεμών decode σταδίων}
Η τρέχουσα υλοποίηση διανυσματοποιεί τα στάδια syndromes, root-finding (FFT), και
error evaluation του \gls{rs} decoder. Το στάδιο Berlekamp-Massey
παραμένει scalar λόγω loop-carried dependencies.

%% --------------------------------------------------------------------------
\label{sec:future-isa}

\paragraph{Εντολή \code{vkeccak.vi} (RISC-V PQC Task Group, υπό τυποποίηση)}
Η ομάδα εργασίας PQC του RISC-V International έχει ξεκινήσει επίσημη διαδικασία
τυποποίησης αφιερωμένης εντολής μετάθεσης Keccak (\code{vkeccak.vi}) υπό τον
κωδικό RVS-2360~\cite{RISCV-PQC-Keccak-Plan2026}. Το δημοσιευμένο πλάνο θέτει
στόχο τελικής τυποποίησης (v1.0) έως τον Δεκέμβριο 2026. 
Η εντολή προτείνεται για επεξεργαστές κλάσης εφαρμογών (application-class) με
vector unit, αξιοποιώντας register group αρκετά μεγάλο ώστε να χωρά την 1600-bit
κατάσταση Keccak ($25\times64$-bit words), λ.χ. LMUL=8 σε \VLEN${}=256$ (group 2048 bits),
αποθηκευμένη σε ομαδοποιημένους διανυσματικούς καταχωρητές. Εφόσον
τυποποιηθεί, θα εξαλείψει τον Keccak πυρήνα που αποτελεί τον επόμενο Amdahl's Law
floor στην τρέχουσα υλοποίηση (\pct{57.0} του Tier~3 \gls{kem} \instret,
\cref{tab:kernel-redistribution}). Σύμφωνα με την πρόταση του PQC Task Group για
την εντολή~\cite{Saarinen2024-vkeccak}, η μετάθεση Keccak-f[1600] (24 γύροι)
συμπυκνώνεται σε κόστος κάτω των \num{100} κύκλων, με επιτάχυνση τάξης
$\geq\!20\times$ σε επίπεδο μετάθεσης και $\sim\!2$–$3\times$ σε επίπεδο σχήματος
για PQC αλγορίθμους. Με μερίδιο Keccak \pct{57.0} στον Tier~3 \gls{hqc} \gls{kem}
\instret, ο νόμος του Amdahl φράσσει το αντίστοιχο scheme-level κέρδος για τον
\gls{hqc} σε $\lesssim\!2.3\times$ (πλήρης εξάλειψη του Keccak). Η ποσοτικοποίηση
αυτή παραμένει εκτίμηση έως την τυποποίηση (RVS-2360) και την εμπειρική επικύρωση
σε silicon.

\paragraph{Configurable-polynomial GF reduction extension}
Η \gls{zvbc} \insn{vclmul.vx} εξυπηρετεί αποτελεσματικά τον πολλαπλασιασμό στο
$\GFpoly$, αλλά η αναγωγή σε μικρά σώματα (\eg $\GF{2^8}$) απαιτεί
πολλαπλές vector εντολές ανά πράξη. Ένα μελλοντικό extension με
configurable reduction polynomial (στο πνεύμα ενεργών προτάσεων για vectorized $\GF{2^8}$ multiply) θα
συμπίεζε την αναγωγή σε μία εντολή, ωφελώντας τόσο τον HQC \gls{rs}
decoder όσο και γενικότερα AES-GCM/GHASH workloads.

\paragraph{Ευρύτερη εφαρμογή σε code-based PQC}
Οι τεχνικές που αναπτύχθηκαν (compile-time specialized Karatsuba, \gls{zvbc} vectorized leaves, scalar
\gls{zbb} Keccak) εφαρμόζονται άμεσα σε άλλα code-based σχήματα που
χρησιμοποιούν πολυωνυμικό πολλαπλασιασμό στο $\GFpoly$, ιδίως στο BIKE, το
οποίο μοιράζεται τον ίδιο αλγεβρικό πυρήνα με τον \gls{hqc}. Η επέκταση της
υλοποίησης σε BIKE θα τεκμηρίωνε τη γενικευσιμότητα της προσέγγισης.

%% --------------------------------------------------------------------------

\vspace{1em}
Συνοψίζοντας, η παρούσα εργασία τεκμηρίωσε ότι οι επικυρωμένες κρυπτογραφικές
επεκτάσεις του RISC-V επαρκούν για να φέρουν μια υλοποίηση σε καθαρό λογισμικό
code-based \gls{pqc} σε επίπεδο απόδοσης συγκρίσιμο με εξειδικευμένες SIMD
αρχιτεκτονικές, χωρίς custom hardware, χωρίς vendor lock-in\ifTODO, και χωρίς
υποχώρηση στην ασφάλεια timing side-channels\fi. Η cryptographic agility που
παρέχει η modular ISA φιλοσοφία του RISC-V αποδεικνύεται πρακτικά εφικτή:
υψηλή απόδοση και αλγοριθμική ευελιξία δεν αποτελούν αντικρουόμενους στόχους.
